% =====================================================================
%  2bat-ud6-1.tex  —  SESSIÓ 1: De la mostra a la conclusió (connexió)
%  (sense preàmbul: s'inclou des de main.tex)
% =====================================================================
\horatitol{Sessió 1 — De la mostra a la conclusió}%
{Connectar amb la normal de la UD5 i obrir la pregunta que mou tota la unitat: com traiem conclusions de tota una població quan només en mesurem una part?}

\subsection{Idea clau}

Arrenquem l'última unitat del curs i de l'etapa. A la UD5 vau aprendre a \textbf{descriure}
el comportament d'un col·lectiu amb les distribucions (binomial i normal): teníeu les
dades de \emph{tothom} i les modelàveu. Avui obrim la pregunta que dona sentit a tot el
que ve: \textbf{i quan no podem mesurar tothom? Quan només tenim una part —una mostra— i
volem treure conclusions sobre tota la població?}

\begin{idea}
\textbf{Població i mostra.}
\begin{itemize}[leftmargin=1.4em, itemsep=2pt]
  \item \textbf{Població:} tot el col·lectiu que es vol conèixer (per exemple, totes les
        persones grans d'un barri).
  \item \textbf{Mostra:} la part que realment s'ha mesurat o enquestat.
  \item \textbf{Inferència:} el \emph{salt} de la mostra a la població, és a dir, treure
        conclusions sobre el tot a partir d'una part.
\end{itemize}
\textbf{La idea central de la unitat:} cap enquesta pregunta a tothom; totes infereixen a
partir d'una mostra. Estudiar com es fa aquest salt —i amb quines garanties i
limitacions— és l'objecte de tota la unitat.
\end{idea}

\subsection{Exemples}

\begin{exemple}[cap estudi pregunta a tothom]
Un titular diu: «el $42\%$ de la població està a favor d'ampliar el servei de menjadors
socials». És impossible que hagin preguntat als milions d'habitants: han enquestat una
\textbf{mostra} (potser $1\,000$ persones) i han \textbf{inferit} el resultat per a tota
la població. Tota la força —i tota la limitació— de l'estudi depèn de com sigui aquesta
mostra.
\end{exemple}

\begin{exemple}[per què la mostra pot representar el tot]
A la UD5 vau veure que moltes magnituds es reparteixen en forma de \textbf{campana}
(normal). Aquesta és, precisament, la raó de fons que fa possible la inferència: si una
mostra es tria bé i és prou gran, tendeix a assemblar-se a la població, i el que hi
observem és una bona \textbf{estimació} del que passa al tot. Però mai és exacta: sempre
hi ha un grau d'incertesa, que aprendrem a expressar.
\end{exemple}

\subsection{Exercicis resolts}

\begin{exresolt}[distingir població i mostra]
Un estudi vol conèixer l'opinió dels usuaris d'un centre cívic (uns $2\,000$) i n'enquesta
$150$. Identifica la població i la mostra.
\solucio
La \textbf{població} és el conjunt de tots els usuaris del centre cívic (uns $2\,000$). La
\textbf{mostra} són les $150$ persones que realment han respost l'enquesta. L'estudi
inferirà l'opinió de tots els $2\,000$ a partir d'aquestes $150$.
\resposta{Població: els $\approx 2\,000$ usuaris; mostra: les $150$ enquestades.}
\end{exresolt}

\begin{exresolt}[per què no es pregunta a tothom]
Per què els estudis solen treballar amb una mostra i no amb tota la població? Dona dues
raons.
\solucio
Perquè mesurar tota la població sol ser \textbf{inviable}: (1) massa \textbf{car} i
\textbf{lent} (preguntar a milions de persones); (2) de vegades \textbf{impossible}
(no es pot localitzar tothom). Una mostra ben triada permet estimar el resultat amb un cost
raonable i una incertesa controlada.
\resposta{Per cost/temps i perquè sovint és impossible mesurar tothom.}
\end{exresolt}

\subsection{Exercicis per fer}

\begin{exercicis}
\begin{exlist}
  \item Un ajuntament vol saber què opinen els joves de la ciutat (uns $30\,000$) i
        n'enquesta $600$. Identifica la població i la mostra.
  \item Digues, en cada cas, si es tracta de la població o d'una mostra:
        \begin{enumerate}[label=\alph*), leftmargin=1.6em, topsep=2pt]
          \item tots els usuaris d'un servei \hfill b) les $80$ persones que han respost
                una enquesta
        \end{enumerate}
  \item Un estudi afirma que «el $60\%$ de les famílies ateses estan satisfetes». Creus
        que han preguntat a totes les famílies o a una mostra? Per què?
  \item Explica amb les teves paraules què vol dir «inferir»: quin és el salt que fa un
        estudi quan passa de la mostra a la població?
  \item \textbf{(Reflexió)} Escriu amb les teves paraules la diferència entre tenir les
        dades de \emph{tothom} i treure conclusions a partir d'una \emph{part}. Quina de
        les dues situacions té més incertesa?
  \item \textbf{(Connexió UD5)} Per què el fet que moltes magnituds es reparteixin «en
        forma de campana» ajuda que una mostra representi bé la població?
\end{exlist}
\end{exercicis}

% ---------------------------------------------------------------------
\fullsolucions{Sessió 1}
\begin{solucions}
\begin{sollist}
  \item Població: tots els joves de la ciutat (uns $30\,000$). Mostra: els $600$
        enquestats.
  \item (a) població (tot el col·lectiu); (b) mostra (la part que ha respost).
  \item A una mostra: preguntar a \emph{totes} les famílies ateses seria molt costós i
        sovint inviable; l'estudi enquesta una part i infereix el resultat.
  \item Resposta oberta. Idea: inferir és treure una conclusió sobre tota la població a
        partir del que s'observa en una mostra; el salt és estendre a tot el col·lectiu el
        que només s'ha mesurat en una part.
  \item Resposta oberta. Idea: amb les dades de tothom hi ha certesa (es coneix el valor
        real); amb una part hi ha \textbf{incertesa}, perquè la mostra pot no coincidir
        exactament amb la població. La segona situació té més incertesa.
  \item Resposta oberta. Idea: si repetíssim el mostreig moltes vegades, els resultats es
        distribuirien en forma de campana al voltant del valor real de la població; per
        això una mostra ben triada i prou gran sol caure a prop d'aquest valor.
\end{sollist}
\end{solucions}
